\section*{Resumen}

Este trabajo presenta la aplicación de la matriz insumo-producto de Leontief a la economía del departamento de La Guajira, Colombia, utilizando tres sectores económicos: Agricultura, Tecnología y Servicios. El objetivo principal es ilustrar el uso y las bondades de esta herramienta para analizar las interacciones entre sectores y el impacto de cambios en la demanda final sobre la producción total requerida.

Se desarrollaron tres ejercicios: (1) cálculo de la matriz de coeficientes técnicos, la matriz inversa de Leontief y los multiplicadores de producción; (2) análisis del impacto de un aumento del 20\% en la demanda final del sector Agricultura sobre la producción de todos los sectores; y (3) comparación de tres escenarios de demanda final (conservador, base y optimista) y su efecto en la producción requerida.

Los resultados muestran que el sector Agricultura tiene un multiplicador de 1.62, el sector Tecnología de 1.67 y el sector Servicios de 1.60. El aumento del 20\% en la demanda de Agricultura genera un incremento del 14.4\% en su propia producción, del 2.0\% en Tecnología y del 1.1\% en Servicios. La comparación de escenarios revela que un escenario optimista (+15\% en demanda) aumenta la producción requerida en un 15.0\% para Agricultura, 15.0\% para Tecnología y 15.0\% para Servicios.

\textbf{Palabras clave:} matriz de Leontief, La Guajira, sectores económicos, multiplicadores, demanda final, producción total.

\clearpage
